% !PRINTED 581
\weberpage{581}
\begin{center}
\textbf{\S\,156.\qquad \gesperrt{Übergang zu den singulären Moduln.}}
\end{center}

Wir nehmen jetzt an, daß $\omega$ die Wurzel einer quadratischen
Gleichung:
\weq{1}{a\omega^2 + b\omega + c = 0}
mit der \gesperrt{negativen Stammdiskriminante}
\weq{2}{\Delta = b^2 - 4ac}
sei.

Ist dann $j(\omega)$ die Invariante, so genügt die Größe $\lambda$ [§ 155,
(2)] der Gleichung 6ten Grades:
\weq{3}{\lambda^6 - 9\cdot16\lambda^4 - [j(\omega)-27\cdot2^8]\lambda^2 + 64[j(\omega)-27\cdot64] = 0,}
deren Koeffizienten ganze algebraische Zahlen sind. Folglich ist
$\lambda$ selbst eine \gesperrt{ganze algebraische Zahl} (Bd. II, § 154, 11.)\footnote{Dieser Umstand ist es, der die Einführung von $\lambda$ an Stelle von $\varkappa$ besonders empfiehlt, weil $\varkappa$ im allgemeinen keine ganze Zahl ist, sondern erst $4\varkappa$.}
und in § 135 haben wir gesehen, daß $\lambda$ Klasseninvariante der
Diskriminante $4\Delta$ oder $16\Delta$ ist.

Wir adjungieren also dem Klassenkörper $\mathfrak{K}(\Delta)$ die Zahl $\lambda$
und erhalten einen Klassenkörper:
\weq{4}{\begin{array}{l}
\mathfrak{L} = \mathfrak{K}(4\Delta), \text{ wenn } \Delta \equiv 0\ (\mathrm{mod}\ 4), \text{ oder } \equiv 1\ (\mathrm{mod}\ 8).\\
\mathfrak{L} = \mathfrak{K}(16\Delta), \qquad \Delta \equiv 5\ (\mathrm{mod}\ 8).
\end{array}}
Ist $h$ die Klassenzahl von $\Delta$, also der Relativgrad von $\mathfrak{K}(\Delta)$
in bezug auf den quadratischen Körper $\Omega = \mathfrak{R}(\sqrt{\Delta})$, und $h'$ der
Relativgrad von $\mathfrak{L}$ in bezug auf $\Omega$, so ist (§ 100, § 123):
\weq{5}{\begin{array}{ll}
h' = h, & \text{wenn } \Delta \equiv 1\ (\mathrm{mod}\ 8),\\
h' = 2h, & \text{,,}\quad \Delta \equiv 0\ (\mathrm{mod}\ 4),\\
h' = 6h, & \text{,,}\quad \Delta \equiv 5\ (\mathrm{mod}\ 8).
\end{array}}
In dem ersten dieser drei Fälle ist also $\mathfrak{L}$ mit $\mathfrak{K}(\Delta)$ identisch.

Durch die Substitution des singulären Moduls $\omega$ gehen die
Koeffizienten $a_1, a_2 \ldots, b_1, b_2, \ldots$, in § 155, (13), (18) in ganze
algebraische Zahlen über, und daraus folgt, Bd. II, § 154, 11.,
daß die Wurzeln von $A(x)$, nämlich:
\weq{6}{\sqrt{\varkappa}\,\mathrm{sn}\,\Omega_{h,h'}, \text{\gesperrt{ganze Zahlen}}}
und die Wurzeln von $B(x)$:
\weq{7}{\sqrt{\varkappa}\,\frac{\mathrm{cn}\,\Omega_{h,h'}}{\mathrm{dn}\,\Omega_{h,h'}}, \text{\gesperrt{Einheiten}}}
sind.
